% =========================================================================
% SECTION 1: INTRODUCTION
% =========================================================================
\section{Introduction}
\label{sec:introduction}

\IEEEPARstart{P}{rimary} mobility aids such as the long white cane and guide dogs reliably detect ground-level hazards within physical reach~\cite{giudice2008blind,elmannai2017sensor,roentgen2008electronic}. However, a standard cane sweep cannot alert blind and visually impaired (BVI) pedestrians to dynamic obstacles approaching from several meters away. Wearable and handheld Electronic Travel Aids (ETAs) attempt to extend this range~\cite{dakopoulos2009wearable}, but most rely on static proximity alarms~\cite{tapu2020wearable,bai2017smart,shoval1998auditory} that sound whenever an obstacle enters a fixed distance boundary.

In shared indoor corridors, static proximity alarms cause two operational failures:
\begin{enumerate}
    \item \textbf{Preceding Pedestrian Alarms}: Walking behind someone at a matched pace triggers continuous alarms despite zero collision risk, inducing alert fatigue~\cite{tapu2020wearable}.
    \item \textbf{Delayed Alerts on Head-On Encounters}: Pedestrians closing head-on at combined walking speeds ($2.0\text{--}2.5$\,m/s) leave under one second to react if an alert triggers at a static 2.5\,m boundary~\cite{lin2021real}.
\end{enumerate}

Active depth sensors (RGB-D, LiDAR) resolve metric scale directly, but their bulk, battery consumption, and cost hinder everyday smartphone deployment~\cite{zhao2020visual}. In contrast, commodity Android smartphones provide integrated cameras, 6-axis IMUs, and multimodal actuators~\cite{faria2010electronic,tapu2020wearable}. However, monocular mobile vision requires overcoming lack of direct depth, handset tilt motion, and frame latency under variable CPU load.

PRISM (Probabilistic Real-time Infrastructure for Safe Mobility) addresses these constraints natively on Android (API 26--34). Operating entirely on-device, PRISM couples monocular ground-plane projection with dynamic IMU tilt compensation, multi-target tracking, analytical Time-to-Collision (TTC) kinematics, and an accelerometer step detector. Navigational guidance is delivered through rate-limited speech, haptics, and 44.1\,kHz stereo spatial audio~\cite{loomis1998navigation,bujacz2016sonification} that conveys safe evasion headings.

\textbf{Contributions:}
\begin{enumerate}
    \item An on-device Android obstacle awareness architecture (API 26--34) with bounded single-frame backpressure ($N=1$).
    \item Monocular ground-plane inverse projection with dynamic IMU gravity compensation and multi-target velocity tracking.
    \item A 9-state hazard lifecycle state machine and multimodal presentation layer using 44.1\,kHz stereo spatial audio to separate obstacle position from heading guidance.
    \item An empirical engineering evaluation comprising a 100-test Android unit suite, 158 controlled encounter playbacks, and an A/B evaluation of sensory feedback modes.
\end{enumerate}
